\section{圆的渐伸线的切线和法线}
\markboth{48. 圆的渐伸线的切线和法线}{48. 圆的渐伸线的切线和法线}

\noindent
\begin{minipage}[c]{0.45\textwidth}
  \centering
  \begin{tikzpicture}[scale=1.00, >=Stealth, line cap=round, line join=round]
    \def\a{0.75}
    \def\phideg{58}
    \pgfmathsetmacro{\phirad}{\phideg*pi/180}
    \coordinate (O) at (0,0);
    \coordinate (A) at (\a, 0);
    \coordinate (B) at ({\a*cos(\phideg)}, {\a*sin(\phideg)});
    \coordinate (P) at ({\a*(cos(\phideg)+\phirad*sin(\phideg))},
                        {\a*(sin(\phideg)-\phirad*cos(\phideg))});
    % C 为从 O 向切线 PT 所作垂线的垂足: OC 方向为 phideg - 90, 长度 a*phirad
    \coordinate (C) at ({\a*\phirad*cos(\phideg-90)}, {\a*\phirad*sin(\phideg-90)});
    % T 在 PT 延长线上
    \coordinate (T) at ($(P)+(\phideg:0.52)$);

    % 坐标轴
    \draw[->] (-0.95, 0) -- (1.72, 0) node[below, font=\scriptsize] {$x$};
    \draw[->] (0, -0.95) -- (0, 1.32) node[left, font=\scriptsize] {$y$};

    % 基圆与渐伸线
    \draw[thick] (O) circle (\a);
    \draw[very thick, curveblue, domain=0:94, samples=60, variable=\t]
      plot ({\a*(cos(\t)+(\t*pi/180)*sin(\t))},
            {\a*(sin(\t)-(\t*pi/180)*cos(\t))});

    % 线段 OB, BP, OC, 切线 CT（phi=58° 使 OB 远离 y 轴、OC 更陡，a 与 A 均位于纯空白区）
    \draw (O) -- node[pos=0.58, above left=-1.5pt, font=\scriptsize] {$a$} (B);
    \draw (B) -- (P);
    \draw (O) -- (C);
    \draw[thick] (C) -- (T);

    % C 处直角标记（清晰远离基圆外缘）
    \coordinate (Cu) at ($(C)+({\phideg}:0.095)$);
    \coordinate (Cv) at ($(C)+({\phideg+90}:0.095)$);
    \draw[thin] (Cu) -- ($(Cu)+(Cv)-(C)$) -- (Cv);

    % 角度 phi 与 alpha
    \draw (0.21, 0) arc (0:\phideg:0.21);
    \node[font=\scriptsize] at (29:0.35) {$\varphi$};

    % 切线 PT 与 x 轴交点 Q
    \pgfmathsetmacro{\Qx}{\a*\phirad/sin(\phideg)}
    \coordinate (Q) at (\Qx, 0);
    \draw ($(Q)+(0:0.18)$) arc (0:\phideg:0.18);
    \node[font=\scriptsize] at ($(Q)+(29:0.32)$) {$\alpha$};

    % 标点与标签（A 置于 x 轴下方、基圆外侧且在 OC 上方的纯空白处）
    \fill (O) circle (1.1pt) node[below left=-1pt, font=\scriptsize] {$O$};
    \fill (A) circle (1.1pt) node[below right=1pt, xshift=1pt, font=\scriptsize] {$A$};
    \fill (B) circle (1.2pt) node[above left=-1pt, font=\scriptsize] {$B$};
    \fill (P) circle (1.2pt) node[right=1pt, font=\scriptsize] {$P$};
    \node[above right=-1pt, font=\scriptsize] at (T) {$T$};
    \fill (C) circle (1.1pt) node[below right=-1pt, font=\scriptsize] {$C$};
  \end{tikzpicture}
  \captionof*{figure}{图\ 48.1}
\end{minipage}%
\hfill
\begin{minipage}[c]{0.52\textwidth}
如图 48.1，设 $P(x,y)$ 是圆的渐伸线上的任意点，$B$ 是基圆上的对应点，则 $PB$ 是基圆的切线，因而 $PB$ 垂直于半径 $OB$。又设
\[
  \angle AOB = \varphi,
\]
则由第 47 节的结果得到
\begin{align*}
  x &= a(\cos\varphi + \varphi\sin\varphi), \\
  y &= a(\sin\varphi - \varphi\cos\varphi).
\end{align*}
用撇号表示对参数 $\varphi$ 的导数，对 $\varphi$ 求导得
\end{minipage}
\begin{align*}
  x' &= a(-\sin\varphi + \sin\varphi + \varphi\cos\varphi) \\
     &= a\varphi\cos\varphi, \\
  y' &= a(\cos\varphi - \cos\varphi + \varphi\sin\varphi) \\
     &= a\varphi\sin\varphi.
\end{align*}
设 $PT$ 是圆的渐伸线在点 $P$ 的切线，$PT$ 与 $x$ 轴的夹角为 $\alpha$，则
\[
  \tg\alpha = \frac{dy}{dx} = \frac{y'\dd\varphi}{x'\dd\varphi} = \frac{a\varphi\sin\varphi}{a\varphi\cos\varphi} = \tg\varphi,
\]
所以 $\alpha=\varphi$，由此推出
\[
  PT \parallel OB.
\]

又因 $PB\perp OB$，所以 $PB\perp PT$，这说明 $PB$ 是圆的渐伸线在 $P$ 点的法线。这样就证明了圆的渐伸线的法线的一个基本性质：

圆的渐伸线在任意点的法线，同时也是基圆的切线。

又因为两平行线间的距离处处相等，所以作切线 $PT$ 的垂线 $OC$（图 48.1），必有 $OC=BP$。而 $BP$ 等于 $\arc{BA}$ 的弧长 $a\varphi$，这里 $a$ 是基圆的半径，所以推出 $OC=a\varphi$。这样就证明了圆的渐伸线的切线的一个重要性质：

从基圆圆心到圆的渐伸线在任意点的切线的距离，与该点的参数值 $\varphi$ 成正比，比例系数等于基圆的半径 $a$。

由于参数值 $\varphi$ 等于切线 $PT$ 与 $x$ 轴的夹角（图 48.1），所以上述性质又可叙述成：

从基圆圆心到圆的渐伸线在任意点的切线的距离，与这切线对 $x$ 轴的倾斜角 $\varphi$ 成正比（设 $x$ 轴通过圆的渐伸线与基圆的交点），比例系数等于基圆半径 $a$。

从上面这些性质推出，当圆的渐伸线的法线 $PB$ 绕基圆圆心匀速旋转时，基圆圆心到对应切线的距离 $OC$ 匀速增大或匀速减小，因此，圆的渐伸线用来作为凸轮的轮廓曲线，可以把匀速旋转变成匀速直线运动。例如图 48.2 所示的是捣矿机中的凸轮机构，图中右下方为凸轮，它带有两个像鹦鹉嘴一样的凸起部分，鹦鹉嘴的凸的一面的轮廓曲线是圆的渐伸线。图中左边的杆是捣杆，杆上固定连接着一个台阶，台阶的底面垂直于杆的中心线，并且与凸轮的侧面相切。当凸轮匀速旋转时，鹦鹉嘴上的圆的渐伸线弧顶起台阶，带动捣杆匀速上升；等到鹦鹉嘴转动到脱离台阶以后，捣杆失去支承力，就因自己的重量而落下，捣击矿石。

\begin{center}
  \begin{minipage}{\textwidth}
    \centering
    \begin{tikzpicture}[scale=0.88, >=Stealth, line cap=round, line join=round]
      % 捣矿机凸轮机构（居中排版防溢出，含上下带斜剖面线固定导套、饱满渐开线凸轮及逆时针旋转箭头）
      % 垂直捣杆: x in [-0.30, 0.00], y in [-1.75, 2.08]
      \draw[thick, fill=white] (-0.30, -1.75) -- (-0.30, 2.08)
        .. controls (-0.22, 2.13) and (-0.12, 2.03) .. (0.00, 2.08)
        -- (0.00, -1.75)
        .. controls (-0.08, -1.70) and (-0.18, -1.80) .. (-0.30, -1.75) -- cycle;
      \draw[thick] (-0.30, 2.08) .. controls (-0.22, 2.03) and (-0.12, 2.13) .. (0.00, 2.08);
      \draw[thick] (-0.30, -1.75) .. controls (-0.22, -1.80) and (-0.12, -1.70) .. (0.00, -1.75);

      % 上固定导套（左右对称两块，带 45 度斜剖面线）
      \foreach \xL/\xR in {-0.58/-0.30, 0.00/0.28} {
        \draw[thick, fill=white] (\xL, 1.56) rectangle (\xR, 1.84);
        \begin{scope}
          \clip (\xL, 1.56) rectangle (\xR, 1.84);
          \foreach \s in {-0.40, -0.32, -0.24, -0.16, -0.08, 0.00, 0.08, 0.16, 0.24, 0.32} {
            \draw[thin, black!75] ({\xL+\s}, 1.54) -- ({\xL+\s+0.34}, 1.88);
          }
        \end{scope}
      }

      % 下固定导套（左右对称两块，带 45 度斜剖面线）
      \foreach \xL/\xR in {-0.58/-0.30, 0.00/0.28} {
        \draw[thick, fill=white] (\xL, -1.46) rectangle (\xR, -1.18);
        \begin{scope}
          \clip (\xL, -1.46) rectangle (\xR, -1.18);
          \foreach \s in {-0.40, -0.32, -0.24, -0.16, -0.08, 0.00, 0.08, 0.16, 0.24, 0.32} {
            \draw[thin, black!75] ({\xL+\s}, -1.48) -- ({\xL+\s+0.34}, -1.14);
          }
        \end{scope}
      }

      % 固定在捣杆上的台阶（凸肩），底面 y = 0.88 水平且与凸轮渐开线严格相切
      \draw[thick, fill=white] (-0.42, 0.96) rectangle (0.36, 1.32);
      \draw[thick, fill=white, rounded corners=1.5pt] (-0.48, 0.88) rectangle (0.44, 0.97);
      \draw[thick, fill=white, rounded corners=1.5pt] (-0.48, 1.31) rectangle (0.44, 1.40);
      % 捣杆穿过台阶内部虚线
      \draw[densely dashed] (-0.30, 0.88) -- (-0.30, 1.40);
      \draw[densely dashed] (0.00, 0.88) -- (0.00, 1.40);

      % 凸轮中心 O_c = (0.56, 0)
      \coordinate (Oc) at (0.56, 0);
      \def\rbcam{0.30}
      \def\thcam{12}
      \def\tcamend{176}
      \pgfmathsetmacro{\tipx}{0.56 + \rbcam*(cos(\thcam+\tcamend)+(\tcamend*pi/180)*sin(\thcam+\tcamend))}
      \pgfmathsetmacro{\tipy}{\rbcam*(sin(\thcam+\tcamend)-(\tcamend*pi/180)*cos(\thcam+\tcamend))}

      % 上凸起（厚实饱满的渐开线“鹦鹉嘴”，在 t=168 deg 即 (0.26, 0.88) 处水平顶住台阶）
      \draw[thick, fill=curveblue!10]
        (Oc) --
        plot[domain=44:\tcamend, samples=44, variable=\t]
          ({0.56 + \rbcam*(cos(\thcam+\t)+(\t*pi/180)*sin(\thcam+\t))},
           {\rbcam*(sin(\thcam+\t)-(\t*pi/180)*cos(\thcam+\t))})
        .. controls (0.26, 0.68) and (0.32, 0.48) .. ($(Oc)+(138:0.40)$) -- cycle;
      % 突出描摹外凸渐开线工作面
      \draw[very thick, curveblue, domain=56:\tcamend, samples=44, variable=\t]
        plot ({0.56 + \rbcam*(cos(\thcam+\t)+(\t*pi/180)*sin(\thcam+\t))},
              {\rbcam*(sin(\thcam+\t)-(\t*pi/180)*cos(\thcam+\t))});

      % 下凸起（关于 Oc 严格中心对称 180 度）
      \draw[thick, fill=curveblue!10]
        (Oc) --
        plot[domain=44:\tcamend, samples=44, variable=\t]
          ({0.56 - \rbcam*(cos(\thcam+\t)+(\t*pi/180)*sin(\thcam+\t))},
           {-\rbcam*(sin(\thcam+\t)-(\t*pi/180)*cos(\thcam+\t))})
        .. controls (0.86, -0.68) and (0.80, -0.48) .. ($(Oc)+(318:0.40)$) -- cycle;
      \draw[very thick, curveblue, domain=56:\tcamend, samples=44, variable=\t]
        plot ({0.56 - \rbcam*(cos(\thcam+\t)+(\t*pi/180)*sin(\thcam+\t))},
              {-\rbcam*(sin(\thcam+\t)-(\t*pi/180)*cos(\thcam+\t))});

      % 凸轮圆毂与轴孔
      \draw[thick, fill=white] (Oc) circle (0.40);
      \draw[thick] (Oc) circle (0.20);
      % 中心十字线
      \draw[densely dash dot, thin] (0.04, 0) -- (1.10, 0);
      \draw[densely dash dot, thin] (0.56, -0.56) -- (0.56, 0.56);

      % 凸轮逆时针旋转方向弧形箭头
      \draw[->, thick] ($(Oc)+(10:0.84)$) arc (10:62:0.84);
    \end{tikzpicture}
    \par\vspace{4pt}{\footnotesize 图\ 48.2}
  \end{minipage}
\end{center}
